%%%%%%%%%%%%%%%%%%%%%%%%%%%%%%%%%%%%%%%%%%%%%%%%%%%%%%%%%%%%%%%%%%%%%%%%%%%%%%%%%%%%%%%%%%%%%%%%
%                                                                                              %
%                             Definicao para a classe Artigo                                   %
%                                                                                              %
%%%%%%%%%%%%%%%%%%%%%%%%%%%%%%%%%%%%%%%%%%%%%%%%%%%%%%%%%%%%%%%%%%%%%%%%%%%%%%%%%%%%%%%%%%%%%%%%

\documentclass[portugues, brazil, a4paper,12pt]{article}
\bibliographystyle{plain}

%%%%%%%%%%%%%%%%%%%%%%%%%%%%%%%%%%%%%%%%%%%%%%%%%%%%%%%%%%%%%%%%%%%%%%%%%%%%%%%%%%%%%%%%%%%%%%%%
%                                                                                              %
%                       Pacotes a utilizar na compilacao do documento                          %
%                                                                                              %
%%%%%%%%%%%%%%%%%%%%%%%%%%%%%%%%%%%%%%%%%%%%%%%%%%%%%%%%%%%%%%%%%%%%%%%%%%%%%%%%%%%%%%%%%%%%%%%%

\usepackage[brazil]{babel}
\usepackage{graphicx}
\usepackage{geometry}
\usepackage[utf8]{inputenc}
\usepackage[T1]{fontenc}
\usepackage{algorithm}
\usepackage{algorithmic}

%%%%%%%%%%%%%%%%%%%%%%%%%%%%%%%%%%%%%%%%%%%%%%%%%%%%%%%%%%%%%%%%%%%%%%%%%%%%%%%%%%%%%%%%%%%%%%%%
%                                                                                              %
%                       Configuracao dos pacotes utilizados no doc.                            %
%                                                                                              %
%%%%%%%%%%%%%%%%%%%%%%%%%%%%%%%%%%%%%%%%%%%%%%%%%%%%%%%%%%%%%%%%%%%%%%%%%%%%%%%%%%%%%%%%%%%%%%%%

\geometry{a4paper,left=3cm,right=3cm,top=2.5cm,bottom=2.93cm}


%%%%%%%%%%%%%%%%%%%%%%%%%%%%%%%%%%%%%%%%%%%%%%%%%%%%%%%%%%%%%%%%%%%%%%%%%%%%%%%%%%%%%%%%%%%%%%%%
%                                                                                              %
%                             Capa do relatorio tecnico                                        %
%                                                                                              %
%%%%%%%%%%%%%%%%%%%%%%%%%%%%%%%%%%%%%%%%%%%%%%%%%%%%%%%%%%%%%%%%%%%%%%%%%%%%%%%%%%%%%%%%%%%%%%%%

\begin{document}

\begin{titlepage}

  \vfill

\begin{figure}[H]
  \centering
  \includegraphics[scale=0.45]{logo.jpg}
\end{figure}

  \vfill

  \begin{center}
    \begin{Large}
      \textbf{INSTITUTO FEDERAL DE EDUCAÇÃO, CIÊNCIA E TECNOLOGIA DE MINAS GERAIS}
    \end{Large}
  \end{center}

  \begin{center}
    \begin{large}
      \textbf{Bacharelado em Ciência da Computação e\\ Engenharia Elétrica} \\[1.4cm] 
    \end{large}
  \end{center}

  \vfill

  \begin{center}
    \begin{large}
      \textbf{Iniciação Científica: \\ Projeto e desenvolvimento de um carro robô controlado por smartphone, utilizando a plataforma Amarino \\ (Arduino e Google Android)} \\
      Relatório e Manual do Módulo \textit{Bluetooth}
    \end{large}
  \end{center}

  \vfill

  \begin{center}
    \begin{large}
      Alunos: \\
        João Paulo Fernandes de Cerqueira César - joaopaulofcc@gmail.com \\
		Rodolfo Labiapari Mansur Guimarães - rodolfolabiapari@gmail.com \\
		Tarlei Almeida - tarleialmeida@hotmail.com
    \end{large}
  \end{center}

\vfill

  \begin{center}
    \begin{large}
      Professores Orientadores: Prof. Dr. Otávio de Souza Martins Gomes e \\ Prof. Dr. Rafael Vinícius Tayette da Nobrega
    \end{large}
  \end{center}

\vfill

  \begin{center}
    \begin{large}
      Formiga - MG \\
      \today \\
    \end{large}
  \end{center}

\clearpage
\tableofcontents 
\end{titlepage}

%%%%%%%%%%%%%%%%%%%%%%%%%%%%%%%%%%%%%%%%%%%%%%%%%%%%%%%%%%%%%%%%%%%%%%%%%%%%%%%%%%%%%%%%%%%%%%%%
%                                                                                              %
%                               Nova Seção do Documento                                        %
%                                                                                              %
%%%%%%%%%%%%%%%%%%%%%%%%%%%%%%%%%%%%%%%%%%%%%%%%%%%%%%%%%%%%%%%%%%%%%%%%%%%%%%%%%%%%%%%%%%%%%%%%

\newpage
\section{Introdução}





%%%%%%%%%%%%%%%%%%%%%%%%%%%%%%%%%%%%%%%%%%%%%%%%%%%%%%%%%%%%%%%%%%%%%%%%%%%%%%%%%%%%%%%%%%%%%%%%
%                                                                                              %
%                               Nova Seção do Documento                                        %
%                                                                                              %
%%%%%%%%%%%%%%%%%%%%%%%%%%%%%%%%%%%%%%%%%%%%%%%%%%%%%%%%%%%%%%%%%%%%%%%%%%%%%%%%%%%%%%%%%%%%%%%%

\newpage
%%%%%%%%%%% Example %%%%%%%%%%%%%%%%%%%%%%%%%%
\section{Bibliografia}
\begin{thebibliography}{100} % 100 is a random guess of the total number of 

%references
\bibitem{robo} roboToSH. Acessado 22/10/2013. Disponível em: http://robotosh.blogspot.jp/2012/07/arduino-jy-mcu-\textit{Bluetooth}.html.

\bibitem{wiki} Arduino - \textit{Bluetooth} HC-05. Acessado 19/10/2013. Disponível em: http://www.tuxti.com.br/wiki/index.php/Arduino\_-\_Bluetooth\_HC-05

\bibitem{seguranca} See why \textit{Bluetooth} technology is simple, secure, everywhere. Acessado 19/10/2013. Disponível em: http://www.bluetooth.com/Pages/Simple-Secure-Everywhere.aspx
\end{thebibliography}

\end{document}